
\documentclass[11pt]{article}
\usepackage[margin=1in]{geometry}
\usepackage{longtable}
\usepackage{array}
\usepackage{hyperref}
\usepackage{amsmath}
\usepackage{amssymb}
\usepackage{booktabs}
\usepackage{enumitem}
\usepackage{microtype}

\title{Finite-Sample SLT Analysis of a Tiny Shakespeare Transformer}
\author{OpenClaw / RelaLeap project notes for Benjamin Goertzel}
\date{2026-07-10}

\begin{document}
\maketitle

\begin{abstract}
This report summarizes two GPU runs of a finite-sample singular learning theory
(SLT) estimator on a small character-level transformer trained on Tiny
Shakespeare. The estimator is a WBIC/SGLD proxy: it is designed to give useful
finite-sample evidence about local learning complexity, but it is not an exact
RLCT oracle. Both runs completed on CUDA. A known product-singularity
calibration benchmark passed in both cases, with the larger run giving
\(\hat\lambda = 0.472 \pm 0.026\) against the analytic target \(\lambda=0.5\).
On the transformer itself, the largest estimated learning coefficients occur in
feed-forward weight matrices, followed by attention projection weights, then
embeddings, while biases and several layer-normalization parameters are near
regular or diagnostic-only. The main interpretation is that the singular
structure measured by this estimator is concentrated in high-dimensional weight
matrices, especially the feed-forward sublayers, not in scalar bias-like
parameters. This is evidence that the estimator pipeline is operational and
qualitatively sensible, not yet a final measurement of exact transformer RLCTs.
\end{abstract}

\section{Purpose and provenance}

The goal of these runs was to test whether the RelaLeap SLT estimator can
produce meaningful block-wise finite-sample learning-complexity measurements on
a real small language-model task, after earlier RelaLeap work used weaker
handcrafted or frozen-parameter proxies. The measured model is a tiny
character-level transformer trained on the Tiny Shakespeare corpus.

\begin{itemize}[leftmargin=*]
\item Small run: \texttt{slt\_validation\_run1}, Runpod A100-SXM4-80GB, command
\texttt{PYTHONPATH=src python3 run\_slt\_gpu.py --device cuda --epochs 3 --n-chains 4 --n-samples 200}.
\item Large run: \texttt{slt\_validation\_large\_run1}, Runpod A100-SXM4-80GB,
command \texttt{PYTHONPATH=src python3 run\_slt\_large.py --device cuda --epochs 3 --n-chains 8 --n-samples 1000}.
\item Source commit for the first run was recorded as RelaLeap branch
\texttt{agent/slt-integration}, commit \texttt{6dca8eb}. The second run used the
same estimator family and preserved its report artifacts locally.
\item Local evidence files: \texttt{projects/relaleap/experiments/slt\_validation\_run1/}
and \texttt{projects/relaleap/experiments/slt\_validation\_large\_run1/}.
\end{itemize}

\section{Short SLT background for non-specialists}

Classical statistical learning theory often assumes that a model is
\emph{regular}: different parameter settings correspond cleanly to different
functions, the Fisher information is non-singular, and local loss surfaces look
approximately quadratic. Neural networks violate these assumptions. They have
symmetries, redundant units, flat directions, dead or unused parameters,
permutation equivalences, and many parameter settings that describe the same or
nearly the same predictor.

Singular learning theory studies this more realistic case. Instead of asking
only how many parameters a model has, SLT asks how the Bayesian evidence or free
energy scales near the set of good predictors. The central quantity is the
\emph{real log canonical threshold} (RLCT), usually written \(\lambda\). In this
report, every \(\hat\lambda\) is an empirical finite-sample estimate or proxy,
not an exact mathematical RLCT.

Very roughly:
\begin{itemize}[leftmargin=*]
\item Larger \(\lambda\) means the parameter block contributes more local
learning complexity under the estimator.
\item Smaller \(\lambda\) means the block is simpler, flatter, more redundant,
or less active around the trained solution.
\item A zero or near-zero value can indicate an effectively flat or inactive
parameter direction under the measured sampling setup.
\item A negative estimate is not physically interpreted as a negative RLCT.
Here it is treated as a diagnostic failure or near-regular/noisy finite-sample
artifact, and the report marks that block as \texttt{diagnostic\_only} rather
than clipping it to zero.
\end{itemize}

\section{What was measured}

The report files contain two levels of measurement.

\subsection{Calibration measurement}

Before trusting the transformer block measurements, the same estimator was run
on a known analytic singularity benchmark, \texttt{product\_singularity\_ab2}.
For this benchmark, the target learning coefficient is \(\lambda=0.5\). Passing
this benchmark does not prove all transformer estimates are correct, but it is a
sanity check that the WBIC/SGLD scaling and energy convention are not completely
wrong.

\subsection{Block-wise transformer measurements}

The transformer was split into 37 parameter blocks: token and position
embeddings, attention projection weights and biases, layer-normalization weights
and biases, feed-forward network weights and biases, and the output projection.
For each block, the estimator reports:

\begin{description}[leftmargin=*]
\item[\(\hat\lambda\)] The estimated finite-sample learning coefficient for that
block. Interpret this as local effective complexity, not just parameter count.
\item[SE] A standard-error estimate for \(\hat\lambda\), reflecting Monte Carlo
uncertainty in the estimator.
\item[Number of parameters] The raw number of trainable scalar parameters in the
block. This is useful context but is not the same thing as learning complexity.
A large block can have low effective complexity if it is redundant or inactive.
\item[Energy ESS per chain] Effective sample size of the sampled energy trace.
Low ESS means the SGLD chain has high autocorrelation and the estimate should be
interpreted cautiously.
\item[Energy Rhat] A between-chain convergence diagnostic. Values near 1 are
best. Values much above 1 mean the chains have not mixed ideally. This report
uses such blocks as informative but finite-sample-limited evidence.
\item[Status] \texttt{calibrated} means the estimator accepted the block under
its current checks. \texttt{diagnostic\_only} means the value should not be used
as a positive SLT claim, usually because the finite-sample estimate was negative
or otherwise not physically interpretable as an RLCT proxy.
\end{description}

\section{Run-level results}

\begin{table}[h]
\centering
\begin{tabular}{lrrrrrr}
\toprule
Run & chains & samples & seconds & blocks & diagnostic & calibration \\
\midrule
Small & 4 & 200 & 401.2 & 37 & 9 & 0.546 $\pm$ 0.033 \\
Large & 8 & 1000 & 3873.4 & 37 & 6 & 0.472 $\pm$ 0.026 \\
\bottomrule
\end{tabular}
\caption{Both SLT validation runs completed. ``Diagnostic'' counts are blocks
marked \texttt{diagnostic\_only}. Calibration is the product-singularity
benchmark estimate, whose analytic target is 0.5.}
\end{table}

The larger run is the better source for interpretation because it used twice as
many chains and five times as many samples per chain. Its calibration estimate
was \(0.472 \pm 0.026\), close to the target 0.5, with calibration Rhat 1.024
and no detected endpoint trend. That is the strongest evidence that this
estimator configuration is approximately sane at this scale.

\section{Conceptual interpretation of the main measurements}

\subsection{Calibration: did the estimator know what a known singularity looks like?}

The calibration benchmark asks the estimator to recover a known SLT value before
it is applied to the transformer. The large run gives \(\hat\lambda=0.472\),
which is within about one standard error of the target 0.5. This is encouraging.
It suggests that the estimator's use of total negative log likelihood, WBIC
scaling, and SGLD sampling is at least directionally correct.

However, this is only one benchmark. It does not prove exactness on neural
networks. It only reduces the risk that we are measuring a meaningless artifact
of the code.

\subsection{Lambda-hat: local learning complexity, not raw size}

A block with more parameters often has a larger \(\hat\lambda\), but the relation
is not one-to-one. The output projection has 4160 parameters and
\(\hat\lambda=47.1\) in the large run, while the token embedding has the same
number of parameters but \(\hat\lambda=0.78\). This means the estimator sees the
output projection as much more locally active or constrained by the task than
the token embedding.

Similarly, attention and feed-forward matrices have the same raw matrix sizes in
some cases, but their estimated complexities differ substantially. This is the
point of an SLT-style measurement: it tries to describe local functional
complexity around the learned solution, not merely count knobs.

\subsection{Standard error: Monte Carlo uncertainty}

The SE values in the large run are small compared with the largest \(\hat\lambda\)
values, especially for the major weight matrices. For example, layer0
feed-forward-up has \(126.1 \pm 0.6\). This does not mean the exact RLCT is known
to that precision. It means that within this finite-sample estimator and its
SGLD traces, the fitted slope is stable. Systematic error from finite data,
block isolation, prior choice, SGLD temperature, or non-convergence may still be
larger than the reported SE.

\subsection{ESS and Rhat: sampler health}

ESS and Rhat say how much to trust the sampling. The calibration benchmark has
good Rhat near 1.02. Many transformer blocks have Rhat between about 1.3 and
3.3, and low energy ESS per chain. That is not ideal. The results are therefore
best viewed as a first calibrated finite-sample map of where complexity appears
to live, not as final high-precision SLT geometry.

\subsection{Diagnostic-only blocks: do not clip bad estimates into good news}

Some bias and normalization blocks have negative \(\hat\lambda\) estimates. A
true RLCT is not negative. Rather than hiding this by clipping to zero, the code
marks such measurements \texttt{diagnostic\_only}. Conceptually, these blocks are
not carrying strong positive complexity evidence under the current estimator.
They may be flat, weakly used, too noisy for the current sampler, or poorly
identified in the block-wise setup.

\section{Large-run block summaries}

\begin{table}[h]
\centering
\begin{tabular}{lrrrrr}
\toprule
Class & blocks & min lambda & mean lambda & max lambda & calibrated \\
\midrule
embedding & 2 & 0.780 & 3.425 & 6.069 & 2/2 \\
attention & 16 & -0.121 & 7.665 & 26.357 & 14/16 \\
feed-forward & 8 & -0.227 & 42.920 & 126.112 & 7/8 \\
normalization & 10 & -0.826 & 0.358 & 1.486 & 7/10 \\
output projection & 1 & 47.138 & 47.138 & 47.138 & 1/1 \\
\bottomrule
\end{tabular}
\caption{Large-run summary by architectural class. Feed-forward and attention
weight blocks dominate the measured local learning complexity.}
\end{table}

\begin{table}[h]
\centering
\begin{tabular}{lrrrr}
\toprule
Block & lambda-hat & SE & parameters & Rhat \\
\midrule
layer0.ffn.up.weight & 126.112 & 0.607 & 16384 & 1.96 \\
layer0.ffn.down.weight & 104.103 & 0.465 & 16384 & 1.94 \\
layer1.ffn.up.weight & 65.923 & 0.330 & 16384 & 2.01 \\
output\_projection.weight & 47.138 & 0.265 & 4160 & 1.67 \\
layer1.ffn.down.weight & 46.311 & 0.218 & 16384 & 1.82 \\
layer1.attention.q.weight & 26.357 & 0.125 & 4096 & 1.86 \\
layer0.attention.q.weight & 24.071 & 0.111 & 4096 & 1.83 \\
layer0.attention.o.weight & 20.139 & 0.107 & 4096 & 1.87 \\
layer1.attention.k.weight & 16.163 & 0.092 & 4096 & 1.72 \\
layer0.attention.k.weight & 12.916 & 0.071 & 4096 & 2.03 \\
\bottomrule
\end{tabular}
\caption{Ten largest calibrated block-wise estimates in the large run.}
\end{table}

\section{Scientific interpretation}

The main structural result is:

\begin{quote}
The measured singular structure of this tiny transformer is concentrated in the
feed-forward and attention weight matrices, especially feed-forward matrices,
while embeddings, biases, and many normalization parameters are lower-complexity
or diagnostic-only under this estimator.
\end{quote}

This is plausible for a trained transformer. Feed-forward layers implement much
of the token-conditional nonlinear transformation. Attention projections govern
contextual routing and mixing. Biases and layer-normalization offsets have far
fewer parameters and often behave more like local shifts than like large
singular submodels. The output projection is also complex, as expected, because
it maps hidden states back to the character distribution.

The result is also useful for RelaLeap. If future residual-layer or
predictive-coding modifications are SLT-guided, the first place to look is not a
single global model score but a block-wise map: which blocks are locally complex,
which blocks are flat or redundant, and whether a proposed intervention reduces
or reorganizes complexity in a way that survives controls.

\section{Limitations and caveats}

\begin{enumerate}[leftmargin=*]
\item These are finite-sample WBIC/SGLD proxies, not exact RLCTs.
\item The transformer block Rhat values are not ideal. More chains, longer runs,
and sample-size sweeps are needed for stronger claims.
\item The block-wise estimates do not automatically add to a correct whole-model
RLCT. Interactions between blocks can matter.
\item The estimator used a Tiny Shakespeare-scale model and corpus. Larger
models may show qualitatively different geometry.
\item The calibration benchmark is encouraging but narrow. A full validation
suite should include regular Gaussian, rank-deficient, product, cusp/crossing,
composition, and RelaLeap-shaped benchmarks.
\item Negative estimates are deliberately left visible and downgraded. This is a
feature, not a bug: the report should fail closed rather than manufacture
positive SLT evidence.
\end{enumerate}

\section{Recommended next measurements}

\begin{enumerate}[leftmargin=*]
\item Repeat the large run with longer chains or more samples to improve ESS and
Rhat for transformer blocks.
\item Run sample-size sweeps in the number of tokens used for WBIC, for example
\(n = 4096, 8192, 16384, 32768\), to test finite-sample scaling.
\item Add whole-model and grouped-module estimates to compare with block-wise
estimates and detect interaction effects.
\item Run matched nulls: shuffled corpus, random-label corpus, and untrained or
partially trained model checkpoints.
\item Compare SLT maps before and after HDPC/ePC or residual-layer interventions.
The scientific question should be whether interventions change the geometry in
a controlled and useful way, not merely whether they change perplexity.
\end{enumerate}

\appendix
\section{Complete large-run block table}

\small
\begin{longtable}{p{0.31\textwidth}p{0.16\textwidth}rrrrp{0.14\textwidth}}
\toprule
Block & Class & lambda-hat & SE & parameters & Rhat & Status \\
\midrule
\endhead
embeddings.token & embedding & 0.780 & 0.019 & 4160 & 1.44 & calibrated \\
embeddings.position & embedding & 6.069 & 0.035 & 8192 & 2.07 & calibrated \\
layer0.attention.q.weight & attention & 24.071 & 0.111 & 4096 & 1.83 & calibrated \\
layer0.attention.q.bias & attention & -0.011 & 0.009 & 64 & 2.16 & diagnostic\_only \\
layer0.attention.k.weight & attention & 12.916 & 0.071 & 4096 & 2.03 & calibrated \\
layer0.attention.k.bias & attention & 0.000 & 0.000 & 64 & 1.00 & calibrated \\
layer0.attention.v.weight & attention & 12.174 & 0.076 & 4096 & 2.09 & calibrated \\
layer0.attention.v.bias & attention & 0.686 & 0.018 & 64 & 1.58 & calibrated \\
layer0.attention.o.weight & attention & 20.139 & 0.107 & 4096 & 1.87 & calibrated \\
layer0.attention.o.bias & attention & 0.615 & 0.032 & 64 & 2.55 & calibrated \\
layer0.ln1.weight & normalization & 1.486 & 0.011 & 64 & 1.85 & calibrated \\
layer0.ln1.bias & normalization & -0.522 & 0.013 & 64 & 1.95 & diagnostic\_only \\
layer0.ln2.weight & normalization & 0.851 & 0.018 & 64 & 2.37 & calibrated \\
layer0.ln2.bias & normalization & 1.456 & 0.015 & 64 & 1.54 & calibrated \\
layer0.ffn.up.weight & feed-forward & 126.112 & 0.607 & 16384 & 1.96 & calibrated \\
layer0.ffn.up.bias & feed-forward & 0.670 & 0.015 & 256 & 1.63 & calibrated \\
layer0.ffn.down.weight & feed-forward & 104.103 & 0.465 & 16384 & 1.94 & calibrated \\
layer0.ffn.down.bias & feed-forward & 0.279 & 0.034 & 64 & 2.92 & calibrated \\
layer1.attention.q.weight & attention & 26.357 & 0.125 & 4096 & 1.86 & calibrated \\
layer1.attention.q.bias & attention & -0.121 & 0.012 & 64 & 2.05 & diagnostic\_only \\
layer1.attention.k.weight & attention & 16.163 & 0.092 & 4096 & 1.72 & calibrated \\
layer1.attention.k.bias & attention & 0.000 & 0.000 & 64 & 1.00 & calibrated \\
layer1.attention.v.weight & attention & 3.926 & 0.039 & 4096 & 1.67 & calibrated \\
layer1.attention.v.bias & attention & 0.557 & 0.015 & 64 & 2.11 & calibrated \\
layer1.attention.o.weight & attention & 5.047 & 0.060 & 4096 & 2.17 & calibrated \\
layer1.attention.o.bias & attention & 0.117 & 0.027 & 64 & 3.17 & calibrated \\
layer1.ln1.weight & normalization & -0.826 & 0.018 & 64 & 1.33 & diagnostic\_only \\
layer1.ln1.bias & normalization & 0.195 & 0.018 & 64 & 1.83 & calibrated \\
layer1.ln2.weight & normalization & 0.139 & 0.012 & 64 & 2.40 & calibrated \\
layer1.ln2.bias & normalization & 0.792 & 0.017 & 64 & 2.44 & calibrated \\
layer1.ffn.up.weight & feed-forward & 65.923 & 0.330 & 16384 & 2.01 & calibrated \\
layer1.ffn.up.bias & feed-forward & 0.186 & 0.009 & 256 & 2.05 & calibrated \\
layer1.ffn.down.weight & feed-forward & 46.311 & 0.218 & 16384 & 1.82 & calibrated \\
layer1.ffn.down.bias & feed-forward & -0.227 & 0.020 & 64 & 3.30 & diagnostic\_only \\
final\_layernorm.weight & normalization & 0.015 & 0.023 & 64 & 1.75 & calibrated \\
final\_layernorm.bias & normalization & -0.006 & 0.040 & 64 & 3.00 & diagnostic\_only \\
output\_projection.weight & output projection & 47.138 & 0.265 & 4160 & 1.67 & calibrated \\
\bottomrule
\end{longtable}
\normalsize

\section{Small-run block table}

\small
\begin{longtable}{p{0.31\textwidth}p{0.16\textwidth}rrrrp{0.14\textwidth}}
\toprule
Block & Class & lambda-hat & SE & parameters & Rhat & Status \\
\midrule
\endhead
embeddings.token & embedding & 0.045 & 0.039 & 4160 & 1.75 & calibrated \\
embeddings.position & embedding & 1.374 & 0.029 & 8192 & 2.15 & calibrated \\
layer0.attention.q.weight & attention & 7.451 & 0.130 & 4096 & 1.98 & calibrated \\
layer0.attention.q.bias & attention & -0.069 & 0.010 & 64 & 1.46 & diagnostic\_only \\
layer0.attention.k.weight & attention & 3.153 & 0.077 & 4096 & 1.80 & calibrated \\
layer0.attention.k.bias & attention & 0.000 & 0.000 & 64 & 1.00 & calibrated \\
layer0.attention.v.weight & attention & 2.166 & 0.067 & 4096 & 1.96 & calibrated \\
layer0.attention.v.bias & attention & 0.500 & 0.052 & 64 & 3.96 & calibrated \\
layer0.attention.o.weight & attention & 3.909 & 0.084 & 4096 & 1.63 & calibrated \\
layer0.attention.o.bias & attention & -0.030 & 0.037 & 64 & 1.35 & diagnostic\_only \\
layer0.ln1.weight & normalization & 0.379 & 0.013 & 64 & 1.72 & calibrated \\
layer0.ln1.bias & normalization & -0.995 & 0.026 & 64 & 1.78 & diagnostic\_only \\
layer0.ln2.weight & normalization & 0.215 & 0.024 & 64 & 1.61 & calibrated \\
layer0.ln2.bias & normalization & 0.062 & 0.033 & 64 & 2.14 & calibrated \\
layer0.ffn.up.weight & feed-forward & 28.381 & 0.547 & 16384 & 1.93 & calibrated \\
layer0.ffn.up.bias & feed-forward & 0.904 & 0.059 & 256 & 2.45 & calibrated \\
layer0.ffn.down.weight & feed-forward & 32.883 & 0.588 & 16384 & 2.11 & calibrated \\
layer0.ffn.down.bias & feed-forward & 0.276 & 0.044 & 64 & 1.69 & calibrated \\
layer1.attention.q.weight & attention & 7.611 & 0.138 & 4096 & 1.65 & calibrated \\
layer1.attention.q.bias & attention & -0.093 & 0.008 & 64 & 1.20 & diagnostic\_only \\
layer1.attention.k.weight & attention & 3.645 & 0.102 & 4096 & 1.88 & calibrated \\
layer1.attention.k.bias & attention & 0.000 & 0.000 & 64 & 1.00 & calibrated \\
layer1.attention.v.weight & attention & 1.522 & 0.061 & 4096 & 1.96 & calibrated \\
layer1.attention.v.bias & attention & 0.501 & 0.020 & 64 & 1.97 & calibrated \\
layer1.attention.o.weight & attention & 1.008 & 0.090 & 4096 & 1.94 & calibrated \\
layer1.attention.o.bias & attention & 0.205 & 0.035 & 64 & 2.09 & calibrated \\
layer1.ln1.weight & normalization & -0.640 & 0.023 & 64 & 1.39 & diagnostic\_only \\
layer1.ln1.bias & normalization & 0.155 & 0.028 & 64 & 1.33 & calibrated \\
layer1.ln2.weight & normalization & -0.496 & 0.019 & 64 & 2.06 & diagnostic\_only \\
layer1.ln2.bias & normalization & 0.471 & 0.025 & 64 & 1.94 & calibrated \\
layer1.ffn.up.weight & feed-forward & 14.862 & 0.300 & 16384 & 2.28 & calibrated \\
layer1.ffn.up.bias & feed-forward & -0.045 & 0.018 & 256 & 3.12 & diagnostic\_only \\
layer1.ffn.down.weight & feed-forward & 14.330 & 0.347 & 16384 & 2.01 & calibrated \\
layer1.ffn.down.bias & feed-forward & -0.105 & 0.022 & 64 & 2.40 & diagnostic\_only \\
final\_layernorm.weight & normalization & -0.811 & 0.022 & 64 & 1.38 & diagnostic\_only \\
final\_layernorm.bias & normalization & 0.344 & 0.056 & 64 & 2.28 & calibrated \\
output\_projection.weight & output projection & 17.161 & 0.352 & 4160 & 1.55 & calibrated \\
\bottomrule
\end{longtable}
\normalsize

\section{Report artifact paths}

\begin{itemize}[leftmargin=*]
\item \texttt{projects/relaleap/experiments/slt\_validation\_run1/artifacts/slt\_tinyshakespeare/slt\_tinyshakespeare\_report.json}
\item \texttt{projects/relaleap/experiments/slt\_validation\_large\_run1/artifacts/slt\_large/slt\_tinyshakespeare\_report.json}
\item \texttt{projects/relaleap/experiments/slt\_validation\_run1/RUN.md}
\item \texttt{projects/relaleap/experiments/slt\_validation\_large\_run1/RUN.md}
\end{itemize}

\end{document}
